\chapter[Image arithmétique finie et factorisation \(468=18\cdot26\)]{Image arithmétique finie et factorisation \(468=18\cdot26\)}
\label{chap:alfabeto-42}
\index{code décimal admissible}
\index{chiffre de contrôle}

L'image en base \(1000\) de l'application finie qui attribue les blocs décimaux est un
sous-ensemble propre des mille blocs décimaux. La loi d'évolution définie dans
\cref{chap:tpk-definicion} et l'arithmétique de \(\mathbb Z/9\mathbb Z\)
déterminent une image de 42 triplets, organisée en six classes de sept. La
même séparation des observables additive et multiplicative factorise ensuite les 468
sorties en 18 signatures additives et 26 signatures multiplicatives.

\section{Primitives et récurrence de l'émetteur fini}

\begin{definition}[Données de départ de \(T_{\min}\)]
\label{pf84:E02-15}
Soit \(D=\{N,E,S,O\}\), avec les vecteurs cardinaux
\[
 v(N)=(-1,0),\quad v(E)=(0,1),\quad
 v(S)=(1,0),\quad v(O)=(0,-1)
 \qquad\text{dans }(\mathbb Z/9\mathbb Z)^2.
\]
La signature temporelle et l'indice multiplicatif sont fixés par
\[
 \operatorname{sgn}_9(t)=
 \begin{cases}
 +1,&t\bmod9\in\{1,4,7\},\\
 -1,&t\bmod9\in\{2,5,8\},\\
 0,&t\bmod9\in\{0,3,6\},
 \end{cases}
\]
et
\[
 \operatorname{ind}_2(1,2,4,8,7,5)=(0,1,2,3,4,5),
\]
avec une valeur nulle en dehors des unités. Une entrée de \(T_{\min}\) est une paire
de germes
\[
 I_9:=\{0,1,\ldots,8\},\qquad
 s_+=(i_+,j_+,d_+),\qquad s_\times=(i_\times,j_\times,d_\times)
 \in I_9^2\times D.
\]
Ces tables, les lecteurs \(L_\Sigma,L_\Pi\), les six blocs de neuf
instants et l'inversion des deux vecteurs lorsque \(27\mid t\) constituent toutes les
données de départ. On n'introduit ni \(\pi\), ni \(e\), ni \(\varphi\), ni aucune autre
valeur cible.
\end{definition}

\begin{algorithm}[Récurrence de 54 pas]
\label{pf84:E02-16}
Au départ, les deux curseurs occupent les positions de leurs germes et portent
les vecteurs \(v(d_+)\), \(v(d_\times)\). Pour
\(t=1,\ldots,54\), dans cet ordre :
\begin{enumerate}
  \item si \(\operatorname{sgn}_9(t)=+1\), on accumule
  \(L_\Sigma(i_+(t),j_+(t))\) et l'on déplace le curseur additif ;
  \item si \(\operatorname{sgn}_9(t)=-1\), on accumule
  \(\operatorname{ind}_2(L_\Pi(i_\times(t),j_\times(t)))\) et l'on déplace le
  curseur multiplicatif ;
  \item si \(\operatorname{sgn}_9(t)=0\), les deux curseurs restent immobiles ;
  \item après cette mise à jour, les deux vecteurs sont inversés si
  \(27\mid t\).
\end{enumerate}
Les accumulateurs sont réinitialisés au début de chaque bloc
\(9(r-1)+1\leq t\leq9r\). Si \(S_r,E_r\) sont leurs valeurs finales, on émet
\[
 U_r=100\bar S_r+10(\bar S_r+\bar E_r)_{10}
       +(3S_r+5E_r+1)_{10},
 \qquad
 \bar S_r=S_r\bmod10,\quad \bar E_r=E_r\bmod10,
\]
et \(w_r=(-U_r)\bmod3\). On obtient ainsi
\[
 T_{\min}(s_+,s_\times)=(U_1,\ldots,U_6)
 \in\mathcal U_6^{\rm cat},
 \qquad
 \Psi(U_6)=(-U_6)\bmod3=:w_6\in\mathbb F_3^6.
\]
\end{algorithm}

\begin{criterion}[Contrôle dual indépendant de la valeur cible]
\label{pf84:E02-17}
La récurrence admet deux formulations indépendantes : factorisation des
canaux et simulation conjointe. Le parcours exhaustif des
\(104\,976\) entrées produit le même sextuplet par les deux voies pour chaque
germe, sans consulter de constante cible. Un écart pour un seul
germe réfuterait l'équivalence.
\end{criterion}

\section{Dynamique bivariée et observables par blocs}

Nous travaillons sur le tore \(\mathbb Z_9^2\). Avec la convention précédente, un
germe complet est
\[
\omega=(p_+,p_\times,h_+,h_\times)
\in(\mathbb Z_9^2)^2\times\{N,E,S,O\}^2,
\]
de sorte que
\[
|\Omega|=9^4\cdot4^2=104\,976.
\]
Dans chaque fenêtre de neuf pas, les temps \(1,4,7\) évaluent l'observable additive,
les temps \(2,5,8\) évaluent l'observable multiplicative et les temps \(3,6,9\)
incrémentent le compteur neutre. Si \(S_m,E_m,T_m\) sont les accumulateurs
de la fenêtre \(m\), l'observable décimale est
\[
d_1\equiv S_m,\qquad
d_2\equiv S_m+E_m,\qquad
d_3\equiv3S_m+5E_m+7T_m\pmod{10},
\]
et le triplet décimal est \(U_m=100d_1+10d_2+d_3\).

\begin{lemma}[Compteur neutre]
Pour toute fenêtre, \(T_m=3\).
\end{lemma}
\begin{proof}
Chaque bloc de neuf temps contient exactement les trois temps neutres
\(3,6,9\). Le compteur est donc incrémenté trois fois.
\end{proof}

\section{Image de l'accumulateur additif : \(\operatorname{Im}S_m=\{6,9,12,15,18,21,24\}\)}

Soit \(s=i+j\pmod9\). L'observable additive visible est
\[
f(s)=\operatorname{dr}_9(s+2)=(2,3,4,5,6,7,8,9,1).
\]
Une translation cardinale modifie \(s\) de \(+1\) ou de \(-1\). Dans une fenêtre, on
lit trois termes consécutifs ; inverser l'orientation inverse seulement leur
ordre. Les neuf sommes cycliques sont
\[
9,12,15,18,21,24,18,12,6.
\]

\begin{lemma}[Image de l'accumulateur additif]
\[
S_m\in\{6,9,12,15,18,21,24\}.
\]
En particulier, \(d_1=S_m\bmod10\) parcourt exactement
\(\{1,2,4,5,6,8,9\}\).
\end{lemma}
\nopagebreak[4]
\begin{proof}
La liste précédente énumère toutes les positions initiales de la trajectoire et les
deux orientations. Leurs résidus décimaux sont au nombre de sept et distincts.\qedhere
\end{proof}

\section{Image de l'accumulateur multiplicatif : \(\operatorname{Im}E_m=\{0,1,3,5,7,9\}\)}

Le groupe des unités \((\mathbb Z/9\mathbb Z)^\times\) est cyclique d'ordre
six. Nous fixons la fonction de coordonnée multiplicative
\[
\ell_\times(1,2,3,4,5,6,7,8,9)=(0,1,0,2,5,0,4,3,0).
\]
Si la ligne ou la colonne parcourue possède un facteur divisible par trois, chaque
évaluation appartient à \(\{3,6,9\}\) et apporte zéro. Si le facteur fixe est une
unité, les sommes de trois termes consécutifs de
\(b\mapsto\ell_\times(\operatorname{dr}_9(ab))\) sont :
\[
\begin{array}{c|c}
a&\text{somme possible}\\
\midrule
1&\{1,3,7,9\}\\
2&\{3,5,9\}\\
4&\{1,5,7\}\\
5&\{1,5,7\}\\
7&\{3,5,9\}\\
8&\{1,3,7,9\}.
\end{array}
\]

\begin{lemma}[Image de l'accumulateur multiplicatif]
\[
E_m\in\{0,1,3,5,7,9\}.
\]
\end{lemma}
\begin{proof}
Le cas non inversible apporte \(0\). La table finie des six unités
contient toutes les possibilités du cas inversible et leur union est
\(\{1,3,5,7,9\}\).
\end{proof}

\section{Congruence décimale de l'alphabet ternaire : unicité du 3e chiffre}

\begin{theorem}[Congruence de contrôle]
Tout triplet \(\overline{d_1d_2d_3}\) émis par le système satisfait
\[
\boxed{d_3\equiv5d_2-2d_1+1\pmod{10}.}
\]
\end{theorem}
\begin{proof}
De \(d_2\equiv S_m+E_m\), on obtient
\(E_m\equiv d_2-d_1\pmod{10}\). Comme \(T_m=3\),
\[
d_3\equiv3d_1+5(d_2-d_1)+21
\equiv5d_2-2d_1+1\pmod{10}.\qedhere
\]
\end{proof}

\begin{theorem}[Image exacte de l'application décimale]
Le système produit exactement \(7\cdot6=42\) triplets décimaux. Regroupés par
\(E=d_2-d_1\pmod{10}\), ce sont :
\[
\begin{array}{c|rrrrrrr}
E&\multicolumn{7}{c}{\text{triplets décimaux}}\\
\midrule
0&114&227&443&556&669&885&998\\
1&129&232&458&561&674&890&903\\
3&149&252&478&581&694&810&923\\
5&169&272&498&501&614&830&943\\
7&189&292&418&521&634&850&963\\
9&109&212&438&541&654&870&983.
\end{array}
\]
\end{theorem}
\begin{proof}
Il existe sept possibilités pour \(d_1\) et six pour \(E\). Pour chaque paire,
\(d_2\equiv d_1+E\) et la congruence de contrôle fixe un unique \(d_3\),
de sorte qu'il existe au plus 42 triplets. La table applique ces deux formules
aux 42 paires et ne contient pas de répétitions. L'énumération des
104\,976 germes retrouve tous les triplets de la table et aucun en dehors
de celle-ci ; la borne est atteinte.
\end{proof}

\section{Factorisation des \(2\) observables}

Soit \(\mathcal S=\mathbb Z_9^2\times\{N,E,S,O\}\), avec
\(|\mathcal S|=324\). À chaque germe additif \(P\in\mathcal S\), on
associe la signature de six fenêtres
\[
s(P)=(S_1,\ldots,S_6)\pmod{10},
\]
et à chaque germe multiplicatif \(T\in\mathcal S\),
\[
e(T)=(E_1,\ldots,E_6)\pmod{10}.
\]
Les deux signatures se combinent coordonnée par coordonnée au moyen de
\[
G(s,e)_k=
100s_k+10(s_k+e_k)_{10}+(3s_k+5e_k+1)_{10}.
\]
L'application \(G\) est injective : à partir d'un triplet, on récupère
\(s_k=d_1\) et \(e_k=d_2-d_1\pmod{10}\), tandis que le troisième chiffre vérifie
l'appartenance à l'image.

\begin{lemma}[Classification exacte du canal additif]
\label{pf84:E02-18}
\[
 |\operatorname{im}s|=18,
\]
et chaque signature additive possède exactement \(18\) préimages.
\end{lemma}
\begin{proof}
La signature dépend uniquement de
\((i+j\bmod9,\varepsilon)\), où \(\varepsilon=+1\) pour \(E,S\) et
\(-1\) pour \(N,O\). Les \(9\cdot2\) paires produisent des signatures distinctes ; chacune
possède neuf positions sur sa diagonale et deux directions, c'est-à-dire
\(18\) préimages. Un nombre différent de signatures ou une classe de taille
différente de \(18\) réfuterait la classification.
\end{proof}

\begin{lemma}[Classification exacte du canal multiplicatif]
\label{pf84:E02-19}
\[
 |\operatorname{im}e|=26.
\]
Il existe \(24\) classes de taille \(8\), une de taille \(24\) et une de taille
\(108\).
\end{lemma}
\begin{proof}
L'énumération des \(9^2\cdot4=324\) germes au moyen de la table explicite
de \(\ell_\times\) donne \(26\) sextuplets
\((E_1,\ldots,E_6)\) et l'histogramme
\[
 24\cdot8+24+108=324.
\]
Une classe supplémentaire, absente ou de taille différente réfuterait le recensement.
\end{proof}

\begin{theorem}[Factorisation exacte et injective du catalogue]
\label{pf84:E02-20}
\[
|\operatorname{im}s|=18,\qquad
|\operatorname{im}e|=26,\qquad
|\operatorname{im}G|=468=18\cdot26.
\]
Chaque signature additive possède 18 préimages. Les signatures multiplicatives présentent
l'histogramme
\[
24\text{ classes de taille }8,\qquad
1\text{ de taille }24,\qquad
1\text{ de taille }108.
\]
\end{theorem}
\begin{proof}
Les deux lemmes précédents donnent les cardinaux des signatures. Dans chaque composante
émise, le chiffre des centaines est \(\bar S_r\), tandis que celui des dizaines est
\(\bar S_r+\bar E_r\bmod10\). Par conséquent,
\[
 \bar S_r=\left\lfloor\frac{U_r}{100}\right\rfloor,\qquad
 \bar E_r=
 \left(\left\lfloor\frac{U_r}{10}\right\rfloor-\bar S_r\right)\bmod10.
\]
Le chiffre des unités vérifie \(3S_r+5E_r+1\bmod10\) ; il n'introduit pas d'autre
identification. Ainsi, \(G\) est injective et son image possède
\(18\cdot26=468\) éléments. Les \(24\) classes multiplicatives de taille
\(8\), combinées aux \(18\) signatures additives, produisent \(432\) sextuplets
de multiplicité \(144\) ; les deux autres classes produisent \(18\) sextuplets de
multiplicité \(432\) et \(18\) de multiplicité \(1944\). En particulier,
\[
 432\cdot144+18\cdot432+18\cdot1944=104\,976.
\]
Deux paires de signatures ayant le même \(U_6\), ou un histogramme qui ne restitue pas le
domaine complet, réfuteraient le théorème.
\end{proof}

La composition complète se factorise donc sous la forme
\[
\mathcal S^2\xrightarrow{s\times e}
\operatorname{im}s\times\operatorname{im}e
\xrightarrow{G}\mathcal U_6
\xrightarrow{\rho_3}\mathcal W_6\subset\mathbb F_3^6.
\]
La dernière projection possède 243 valeurs. La première descente conserve la
séparation somme--produit ; la deuxième conserve 468 états décimaux ; la
troisième élimine la coordonnée de fibre et retient un mot ternaire.

\begin{corollary}[Recensement de la projection observable]
\label{cor:censo-proyeccion-observable}
\label{pf84:E02-21}
L'énumération complète détermine la chaîne de cardinaux
\[
 104\,976\longrightarrow468\longrightarrow243\subset729.
\]
Les \(468\) éléments intermédiaires sont des sextuplets d'entiers compris entre
zéro et \(999\) ; seule leur image finale appartient à \(\mathbb F_3^6\). Les
multiplicités des fibres de la première application satisfont
\[
 432\cdot144+18\cdot432+18\cdot1944=104\,976.
\]
\end{corollary}

\begin{proof}
Le théorème de factorisation donne \(18\cdot26=468\), et l'énumération
exhaustive des \(104\,976\) germes produit exactement les trois
multiplicités indiquées. La réduction composante par composante parcourt les
\(243\) mots enregistrés dans le catalogue ternaire. Les catalogues de
signatures et de routes joints reproduisent les deux énumérations.
\end{proof}

\begin{proposition}[Microfibre quintuple du mot \(010211\)]
\label{pf84:E02-22}
Pour la réduction restreinte
\[
 \Psi:\mathcal U_6^{\rm cat}\longrightarrow\mathbb F_3^6,
 \qquad \Psi(U_6)=(-U_6)\bmod3,
\]
la fibre de \(010211\) se compose exactement des cinq sextuplets suivants :
\begingroup
\scriptsize
\[
\begin{array}{@{}cclc@{}}
\toprule
U_6&\sigma_+&\sigma_\times&\text{multiplicité}\\
\midrule
(501,614,498,169,272,272)&(5,6,4,1,2,2)&(5,5,5,5,5,5)&432\\
(810,923,870,169,272,272)&(8,9,8,1,2,2)&(3,3,9,5,5,5)&144\\
(810,983,810,169,272,272)&(8,9,8,1,2,2)&(3,9,3,5,5,5)&144\\
(870,923,810,169,272,272)&(8,9,8,1,2,2)&(9,3,3,5,5,5)&144\\
(870,923,810,418,674,521)&(8,9,8,4,6,5)&(9,3,3,7,1,7)&144\\
\bottomrule
\end{array}
\]
\endgroup
Ce sont cinq préimages du catalogue \(U_6\), non cinq histoires complètes
TPK. Elles conservent des signatures et des multiplicités distinctes même si elles partagent le même
mot visible.
\end{proposition}

\begin{proof}
Le balayage exhaustif du catalogue de \(468\) sextuplets sélectionne exactement
les cinq lignes de la table. La fonction de récupération des signatures reproduit
les deux sextuplets de chaque ligne et la somme des multiplicités est
\(432+4\cdot144=1008\). L'énumération vérifie en outre que chaque cellule est le
produit cartésien de ses germes de canal. Un cardinal différent de cinq
ou l'identification de deux lignes malgré leurs données enrichies réfuterait
l'énoncé.
\end{proof}
